\subsection{弗罗贝尼乌斯互反律}

设 A 是一个半单环。设 f 是从 A 到半单环 B 的一个同态。设 S 是一个单 A-模，T 是一个单 B-模，且设 D 与 E 分别是 S 与 T 的中心化子。由舒尔引理（VIII, 第 47 页，推论），D 与 E 都是体。用 H 表示从 S 到 \(f_{*}(\mathrm{T})\) 的 A-线性同态所成之集。我们赋予 H 一个 (E,D)-双模结构，其外部合成律为 \((e,u)\mapsto e\circ u\) 与 \((d,u)\mapsto u\circ d\)（对 \(e\in E\)，\(u\in H\)，\(d\in D\)）。

命题 11. —— a) 重数 \([f_{*}(\mathrm{T}):\mathrm{S}]\)（单 A-模 S 在半单 A-模 \(f_{*}(\mathrm{T})\) 中的重数）等于把 H 视为 D 上右向量空间时的维数。

\begin{enumerate}
\def\labelenumi{\alph{enumi})}
\setcounter{enumi}{1}
\tightlist
\item
  重数 \([f^{*}(\mathrm{S}):\mathrm{T}]\) 是有限的，且等于把 H 视为 E 上左向量空间时的维数。
\end{enumerate}

断言 a) 由 VIII, 第 72 页的公式 (11) 得出。

B-模 \(f^{*}(S)\) 是半单且有限生成的，故具有有限长度。由同前所引的公式 (12)，我们有

\[
[ f ^ {*} (\mathrm{S}): \mathrm{T} ] = \dim_ {\mathrm{E}} \operatorname{Hom} _ {\mathrm{B}} (f ^ {*} (\mathrm{S}), \mathrm{T}).\tag{20}
\]

另外，在 II, §5, No.~1, 第 277--278 页的公式 (2) 与注 2 中，我们定义了一个从 \(\mathrm{Hom}_{\mathrm{A}}(\mathrm{S},f_{*}(\mathrm{T}))\) 到 \(\mathrm{Hom}_{\mathrm{B}}(f^{*}(\mathrm{S}),\mathrm{T})\) 的 E-线性双射。于是断言 b) 由公式 (20) 得出。

推论（弗罗贝尼乌斯互反律）. —— 假设 A 与 B 是交换域 K 上为有限维的半单代数，且 f 是 K-线性的。则 K-向量空间 S、T、D、E 与 H 都是有限维的，并且我们有等式

\[
[ f _ {*} (\mathrm{T}): \mathrm{S} ] [ \mathrm{D}: \mathrm{K} ] = [ f ^ {*} (\mathrm{S}): \mathrm{T} ] [ \mathrm{E}: \mathrm{K} ] = [ \mathrm{H}: \mathrm{K} ].\tag{21}
\]

特别地，当 K 是代数闭的时，我们有 D = E = K，且

\[
[ f _ {*} (\mathrm{T}): \mathrm{S} ] = [ f ^ {*} (\mathrm{S}): \mathrm{T} ] = [ \mathrm{H}: \mathrm{K} ].\tag{22}
\]

由于 A-模 S 是单的，它是单生成的，且 D 是 \(\operatorname{Hom}_{\mathrm{K}}(\mathrm{S},\mathrm{S})\) 的一个线性子空间；故 S 与 D 在 K 上都是有限维的。出于同样的理由，T 与 E 在 K 上是有限维的。最后，H 是 \(\operatorname{Hom}_{\mathrm{K}}(\mathrm{S},\mathrm{T})\) 的一个线性子空间；因而它也是有限维的。于是公式 (21) 由命题 11 得出，因为 H 在 K 上的维数等于 \([H:D][D:K]\)，也等于 \([H:E][E:K]\)（II, §1, No.~13, 第 222 页，命题 25）。由 VIII, 第 47 页的定理 1，若 K 是代数闭的，我们有 D = E = K。于是推论的第二部分由第一部分得出。

设 \(\mathrm{A}\) 与 \(\mathrm{B}\) 是交换域 \(\mathrm{K}\) 上为有限维的半单代数，且设 \(f\) 是一个 \(\mathrm{K}\)-代数同态（从 \(\mathrm{A}\) 到 \(\mathrm{B}\)）。

用 \(\mathcal{S}(\mathrm{A})\) 表示单 A-模的类所成之集；对每个 \(\lambda \in \mathcal{S}(\mathrm{A})\)，设 \(\mathrm{S}_{\lambda}\) 是类为 \(\lambda\) 的一个模，\(\mathrm{D}_{\lambda}\) 是它的中心化子。则 \(\mathrm{D}_{\lambda}\) 是 \(\mathrm{K}\) 上的一个有限次数代数；我们把这个次数记作 \(d_{\lambda}\)。我们类似地定义 \(\mathcal{S}(\mathrm{B})\)、\(\mathrm{T}_{\mu}\)、\(\mathrm{E}_{\mu}\) 与 \(e_{\mu}\)（对 \(\mu\) 属于 \(\mathcal{S}(\mathrm{B})\)）。格罗滕迪克群 \(\mathrm{K}_0(\mathrm{A})\) 以族 \(([\mathrm{S}_{\lambda}])_{\lambda \in \mathcal{S}(\mathrm{A})}\) 为基，而 \(\mathrm{K}_0(\mathrm{B})\) 以族 \(([\mathrm{T}_{\mu}])_{\mu \in \mathcal{S}(\mathrm{B})}\) 为基。设 \((a_{\mu \lambda})\) 是 \(f^{*}: \mathrm{K}_{0}(\mathrm{A}) \to \mathrm{K}_{0}(\mathrm{B})\) 关于这些基的矩阵，\((b_{\lambda \mu})\) 是 \(f_{*}: \mathrm{K}_{0}(\mathrm{B}) \to \mathrm{K}_{0}(\mathrm{A})\) 关于这些基的矩阵。按定义，我们有

\[
a _ {\mu \lambda} = [ f ^ {*} (\mathrm{S} _ {\lambda}): \mathrm{T} _ {\mu} ], \qquad b _ {\lambda \mu} = [ f _ {*} (\mathrm{T} _ {\mu}): \mathrm{S} _ {\lambda} ]\tag{23}
\]

（\(\lambda\) 属于 \(\mathcal{S}(\mathrm{A})\)，\(\mu\) 属于 \(\mathcal{S}(\mathrm{B})\)）。我们把向量空间 \(\mathrm{Hom}_{\mathrm{A}}(\mathrm{S}_{\lambda}, f_{*}(\mathrm{T}_{\mu}))\) 在域 K 上的维数记作 \(h_{\lambda \mu}\)。由上面的推论，我们有

\[
h _ {\lambda \mu} = e _ {\mu} a _ {\mu \lambda} = d _ {\lambda} b _ {\lambda \mu}.\tag{24}
\]

当域 K 是代数闭的时，我们有 \(d_{\lambda} = e_{\mu} = 1\)；于是，我们有

\[
a _ {\mu \lambda} = b _ {\lambda \mu} = h _ {\lambda \mu}.\tag{25}
\]

换句话说，\(f_{*}\) 与 \(f^{*}\) 关于 \(\mathrm{K}_0(\mathrm{A})\) 与 \(\mathrm{K}_0(\mathrm{B})\) 的给定基的矩阵互为转置。

